% Edition check: concise; total-page review limit 5, including references.
% Reader task: Define the objects and central question; complete a model that explains the general structure.
% Select one complete thread; remove genuine side branches before drafting.
% In the Introduction state every selected principal conclusion, not only names.
% Review body-to-introduction reverse coverage after changing any result or scope.
% Style: delete trivial transitions; retain useful self-contained repetition;
% use itemize/enumerate; put secondary consequences afterward, as corollaries when useful.
% An introduction consisting of definitions and theorems is valid:
% definitions and theorems are valid without a narrative roadmap.
% Two slots illustrate wiring, not a theorem-count requirement. Remove unused slots.
\documentclass[12pt,reqno]{amsart}
\usepackage{math-review}
\title{[INSERT: Topic]: Foundations and Models}
\author{}
\date{[INSERT: Checked literature cutoff]}
\newcommand{\PrincipalStatement}{%
\placeholder{Specify objects, all necessary hypotheses, quantifiers and ranges.}
\begin{itemize}
\item \placeholder{A necessary condition; omit this list if there is no real list.}
\item \placeholder{A second independent condition, only when needed.}
\end{itemize}
\placeholder{State the actual mathematical conclusion and exact source/status.}
}
\newcommand{\SecondStatement}{%
\placeholder{State another selected principal answer with its own data, necessary conditions and actual conclusion; remove this slot when outside scope.}
}
\begin{document}
\begin{abstract}
\placeholder{Give the question and selected answers or boundary in two or three sentences. The abstract does not replace the main statements below.}
\end{abstract}
\maketitle

\section{Introduction}\label{sec:introduction}
\placeholder{Define the minimum necessary notation. State fixed data, unknowns and branch conditions before using them. Give actual answers with quantifiers, parameter ranges and source/status qualifications. Delete transitions that merely announce what follows.}
\begin{theorem}\label{thm:principal}
\PrincipalStatement
\end{theorem}
\begin{theorem}\label{thm:second}
\SecondStatement
\end{theorem}
\placeholder{Keep a precise remaining question or negative boundary here when it is a principal commitment. Change theorem to proposition, sourced claim or Problem as mathematically appropriate. No invented optional titles.}

\section{Objects and a completed model}\label{sec:mechanism}
\begin{theoremrecall}{thm:principal}
\PrincipalStatement
\end{theoremrecall}
\placeholder{Carry out the chosen model calculation to its conclusion. State exactly which feature survives in the general setting; do not list examples without calculations. A recall is optional if it adds backtracking rather than removing it. Reuse the canonical body whenever a full repetition is retained.}

\section{The second answer and the boundary}\label{sec:boundary}
\begin{theoremrecall}{thm:second}
\SecondStatement
\end{theoremrecall}
\placeholder{Use a mathematical heading. Explain the second answer or boundary without bundling secondary consequences into its main statement. Define a precise sourced Problem only if the selected scope needs one. Do not manufacture a fixed number of Problems.}

% Replace, relocate or delete this structural model according to actual mathematics.
% Do not leave core body results hidden inside the preceding explanatory paragraphs.
\begin{lemma}\label{lem:body-input}
\placeholder{State a necessary supporting input with local hypotheses and a usable output. No fixed lemma quota is imposed.}
\end{lemma}
\begin{proof}[Proof sketch]\label{prf:body-input}
\placeholder{Separate the argument from the statement; give a decisive step and cite any imported estimate precisely. Standard deepens this proof, not the number of wrappers.}
\end{proof}

\templatereferences
\end{document}
